% =====================================================================
%  preamble.tex  —  all styling lives here (Notion-like look)
%  Compile with XeLaTeX or LuaLaTeX (Overleaf: Menu → Compiler)
% =====================================================================

\usepackage[a4paper, margin=2.4cm, bottom=2.6cm]{geometry}
\usepackage{iftex}

% ---------- Font: Inter (bundled in fonts/) --------------------------
\iftutex
  \usepackage{fontspec}
  \setmainfont{Inter}[
    Path           = fonts/,
    Extension      = .otf,
    UprightFont    = *-Regular,
    ItalicFont     = *-Italic,
    BoldFont       = *-SemiBold,
    BoldItalicFont = *-SemiBoldItalic,
    Scale          = 0.95]
  \newfontfamily\headingfont{Inter}[
    Path = fonts/, Extension = .otf,
    UprightFont = *-Bold, ItalicFont = *-BoldItalic]
\else
  % fallback if compiled with pdfLaTeX
  \usepackage[T1]{fontenc}
  \usepackage[scaled]{helvet}
  \renewcommand{\familydefault}{\sfdefault}
  \newcommand{\headingfont}{\bfseries}
\fi

\usepackage[english]{babel}
\usepackage{microtype}
\linespread{1.15}

% ---------- Packages ------------------------------------------------
\usepackage{graphicx}
\usepackage[table]{xcolor}
\usepackage{tabularx, array}
\usepackage{float}
\usepackage{pdfpages}
\usepackage{enumitem}
\usepackage[skip=8pt plus 2pt]{parskip}
\usepackage{fancyhdr}
\usepackage{titlesec}
\usepackage{tocloft}
\usepackage[most]{tcolorbox}
\usepackage[hidelinks]{hyperref}
\usepackage{bookmark}
\usepackage{needspace}
\usepackage{ragged2e}
\AtBeginDocument{\RaggedRight}   % left-aligned like Notion, no justified gaps

% ---------- Colours ---------------------------------------------------
\definecolor{ink}{HTML}{37352F}       % Notion text colour
\definecolor{muted}{HTML}{787774}     % Notion grey
\definecolor{rule}{HTML}{E9E9E7}      % Notion table lines
\definecolor{shade}{HTML}{F7F6F3}     % Notion table header / callout
\definecolor{linkc}{HTML}{37352F}
\definecolor{checkbg}{HTML}{FBF3DB}   % Notion yellow highlight
\color{ink}
\hypersetup{colorlinks=true, linkcolor=ink, urlcolor=muted}

% ---------- Headings --------------------------------------------------
%  \chapter    = part of the portfolio (Seminars, Debates, ...)
%  \section    = one seminar / debate          (1.1)
%  \subsection = About, Notes, Reflection, ... (1.1.1)
% Each seminar/debate starts on a new page (except the first one after a chapter title)
\newcommand{\sectionbreak}{\ifnum\value{section}>0 \clearpage\fi}

\titleformat{\chapter}[block]
  {\headingfont\fontsize{30}{36}\selectfont}
  {\thechapter.}{0.5em}{}
\titlespacing*{\chapter}{0pt}{0pt}{24pt}

\titleformat{\section}
  {\headingfont\fontsize{19}{24}\selectfont}{\thesection.}{0.5em}{}
\titlespacing*{\section}{0pt}{20pt}{8pt}

\titleformat{\subsection}
  {\headingfont\fontsize{14.5}{18}\selectfont}{\thesubsection.}{0.5em}{}
\titlespacing*{\subsection}{0pt}{16pt}{4pt}

\titleformat{\subsubsection}
  {\headingfont\normalsize}{\thesubsubsection.}{0.5em}{}
\titlespacing*{\subsubsection}{0pt}{12pt}{2pt}

\setcounter{secnumdepth}{3}

% ---------- Table of contents ----------------------------------------
\setcounter{tocdepth}{1}   % 0 = parts only, 1 = + each seminar, 2 = + About/Notes/...
\renewcommand{\cfttoctitlefont}{\headingfont\fontsize{30}{36}\selectfont}
\renewcommand{\cftchapfont}{\bfseries}
\renewcommand{\cftchappagefont}{\bfseries}
\renewcommand{\cftchapleader}{\cftdotfill{\cftdotsep}}
\renewcommand{\cftdotsep}{2.5}
\setlength{\cftbeforechapskip}{10pt}

% ---------- Header / footer ------------------------------------------
\pagestyle{fancy}
\fancyhf{}
\renewcommand{\headrulewidth}{0pt}
\fancyfoot[L]{\footnotesize\color{muted}\PortfolioFooter}
\fancyfoot[R]{\footnotesize\color{muted}\thepage}
\fancypagestyle{plain}{\fancyhf{}\renewcommand{\headrulewidth}{0pt}%
  \fancyfoot[L]{\footnotesize\color{muted}\PortfolioFooter}%
  \fancyfoot[R]{\footnotesize\color{muted}\thepage}}

% ---------- Lists (tight, like Notion) -------------------------------
\setlist{itemsep=2pt, topsep=2pt, parsep=0pt}
\setlist[itemize,2]{label=$\circ$}
\setlist[itemize,3]{label=\tiny\raisebox{0.3ex}{\rule{0.5em}{0.5em}}}

% =====================================================================
%  Helper commands
% =====================================================================

% --- \img[width]{file} : an image; grey placeholder if file missing ---
\newcommand{\img}[2][\linewidth]{%
  \IfFileExists{#2}%
    {\includegraphics[width=#1, height=0.42\textheight, keepaspectratio]{#2}}%
    {\fcolorbox{rule}{shade}{\parbox[c][3cm][c]{\dimexpr#1-2\fboxsep-2\fboxrule}%
      {\centering\footnotesize\color{muted}missing image:\\\texttt{\detokenize{#2}}}}}%
}

% --- Image rows ---------------------------------------------------------
%   \begin{imagerow}
%     \rowimg{0.3}{images/a.jpg}{Caption}  \hfill
%     \rowimg{0.3}{images/b.jpg}{}         \hfill
%     \rowimg{0.3}{images/c.jpg}{Caption}
%   \end{imagerow}
\newenvironment{imagerow}
  {\par\medskip\noindent\begin{minipage}{\linewidth}}
  {\end{minipage}\par\medskip}
\newcommand{\rowimg}[3]{%
  \begin{minipage}[t]{#1\linewidth}\centering
    \img{#2}%
    \ifx\relax#3\relax\else\\[2pt]{\footnotesize\color{muted}#3}\fi
  \end{minipage}}

% --- Text next to an image -----------------------------------------------
%   \sidebyside{0.5}{text ...}{images/x.jpg}    (0.5 = width of the text)
\newcommand{\sidebyside}[3]{%
  \par\medskip\noindent
  \begin{minipage}[c]{#1\linewidth}\RaggedRight #2\end{minipage}\hfill
  \begin{minipage}[c]{\dimexpr0.96\linewidth-#1\linewidth\relax}\centering\img{#3}\end{minipage}
  \par\medskip}

% --- Info table (grey first column, like Notion) ------------------------
%   \begin{infotable}
%     Title & ... \\
%     Speaker & ... \\
%   \end{infotable}
\newenvironment{infotable}
  {\par\smallskip\noindent\arrayrulecolor{rule}\renewcommand{\arraystretch}{1.45}%
   \tabularx{\linewidth}{|>{\columncolor{shade}\raggedright\arraybackslash}p{3.3cm}|>{\raggedright\arraybackslash}X|}\hline}
  {\\\hline\endtabularx\par\smallskip}
\newcommand{\row}{\\\hline}   % use \row between rows instead of \\

% --- Callout box (for reflections) ---------------------------------------
\newtcolorbox{callout}[1][]{enhanced, colback=shade, colframe=shade,
  boxrule=0pt, arc=4pt, left=10pt, right=10pt, top=8pt, bottom=8pt, #1}

% --- Highlight something to double-check before submitting -------------
\newcommand{\tocheck}[1]{\colorbox{checkbg}{\parbox{\dimexpr\linewidth-2\fboxsep}{#1}}}

% --- Insert a signed / scanned PDF as full pages ------------------------
%   \attachpdf{TOC title}{pdfs/file.pdf}{toc level: chapter|section|subsection}
%   The pages get a page number and a TOC entry; a placeholder page
%   is printed if the file is not there yet.
\newcommand{\attachpdf}[3]{%
  \clearpage\phantomsection
  \addcontentsline{toc}{#3}{#1}%
  \IfFileExists{#2}%
    {\includepdf[pages=-, scale=0.9, frame, pagecommand={\thispagestyle{fancy}}]{#2}}%
    {\thispagestyle{fancy}\vspace*{\fill}
     \begin{center}
     \fcolorbox{rule}{shade}{\parbox[c][9cm][c]{0.8\linewidth}{\centering
       {\large\headingfont #1}\\[10pt]
       \color{muted}Placeholder — upload the PDF to\\[2pt]\texttt{\detokenize{#2}}}}
     \end{center}
     \vspace*{\fill}\clearpage}%
}

% --- Unnumbered chapter that still appears in the TOC -------------------
\newcommand{\plainchapter}[1]{\chapter*{#1}\addcontentsline{toc}{chapter}{#1}\markboth{#1}{}}
